\documentclass[11pt,a4paper]{article}

% Document autonome, compilable sur Overleaf avec pdfLaTeX.
\usepackage[utf8]{inputenc}
\usepackage[T1]{fontenc}
\usepackage[french]{babel}
\usepackage{lmodern}
\usepackage{microtype}
\usepackage[a4paper,margin=2.4cm,headheight=15pt]{geometry}
\usepackage{xcolor}
\usepackage{tabularx}
\usepackage{array}
\usepackage{enumitem}
\usepackage{titlesec}
\usepackage{fancyhdr}
\usepackage{tikz}
\usetikzlibrary{arrows.meta,positioning}
\usepackage[most]{tcolorbox}
\usepackage{hyperref}

\definecolor{Navy}{HTML}{14283F}
\definecolor{Blue}{HTML}{245B78}
\definecolor{Teal}{HTML}{168C8C}
\definecolor{PaleTeal}{HTML}{E8F5F3}
\definecolor{PaleBlue}{HTML}{EDF3F8}
\definecolor{Ink}{HTML}{263442}
\definecolor{Muted}{HTML}{607080}
\definecolor{Gold}{HTML}{D9A441}
\definecolor{PaleGold}{HTML}{FBF5E8}
\definecolor{Line}{HTML}{D5DEE6}

\color{Ink}
\hypersetup{
  colorlinks=true,
  linkcolor=Blue,
  urlcolor=Teal,
  pdftitle={Proposition de sujet de projet d'expertise — Sécurité applicative et DevSecOps},
  pdfauthor={Badr Laklach et Taha Benahmadi},
  pdfsubject={Cybersécurité, DevOps, DevSecOps et agents d'intelligence artificielle}
}

\titleformat{\section}
  {\Large\bfseries\color{Navy}}{\thesection}{0.7em}{}
  [\vspace{0.2em}{\color{Teal}\titlerule[1.1pt]}]
\titleformat{\subsection}
  {\large\bfseries\color{Blue}}{\thesubsection}{0.65em}{}
\titlespacing*{\section}{0pt}{1.35em}{0.65em}
\titlespacing*{\subsection}{0pt}{0.9em}{0.35em}

\pagestyle{fancy}
\fancyhf{}
\fancyhead[L]{\small\color{Muted}Projet d'expertise — proposition}
\fancyhead[R]{\small\color{Muted}GI--ILSI}
\fancyfoot[L]{\small\color{Muted}Badr Laklach et Taha Benahmadi}
\fancyfoot[R]{\small\color{Muted}\thepage}
\renewcommand{\headrulewidth}{0.4pt}
\renewcommand{\headrule}{\hbox to\headwidth{\color{Line}\leaders\hrule height \headrulewidth\hfill}}

\setlist[itemize]{leftmargin=1.35em,itemsep=0.25em,topsep=0.35em,label=\textcolor{Teal}{\textbullet}}
\setlength{\parskip}{0.35em}
\setlength{\parindent}{0pt}
\setlength{\emergencystretch}{1.5em}
\clubpenalty=10000
\widowpenalty=10000
\displaywidowpenalty=10000
\newcolumntype{Y}{>{\raggedright\arraybackslash}X}

\newtcolorbox{summarybox}{
  enhanced,breakable,colback=PaleBlue,colframe=Blue!35,
  boxrule=0.5pt,arc=2mm,left=4mm,right=4mm,top=3mm,bottom=3mm,
  borderline west={3pt}{0pt}{Teal}
}
\newtcolorbox{questionbox}[1][Problématique]{
  enhanced,breakable,colback=PaleTeal,colframe=Teal,
  boxrule=0.6pt,arc=2mm,left=4mm,right=4mm,top=3mm,bottom=3mm,
  title=#1,fonttitle=\bfseries\small,coltitle=white,colbacktitle=Teal,
  attach boxed title to top left={yshift=-2mm,xshift=4mm},
  boxed title style={arc=1mm,boxrule=0pt}
}

\begin{document}

% Couverture
\begin{titlepage}
\thispagestyle{empty}
\begin{tikzpicture}[remember picture,overlay]
  \fill[Navy] (current page.north west) rectangle ([yshift=-10.2cm]current page.north east);
  \fill[Teal] (current page.north west) rectangle ([yshift=-10.2cm,xshift=0.38cm]current page.north west);
  \fill[Gold] ([yshift=-10.2cm]current page.north west) rectangle ([yshift=-10.32cm]current page.north east);
  \fill[white,opacity=0.05] ([xshift=-2.0cm,yshift=-2.4cm]current page.north east) circle (3.0cm);
  \fill[white,opacity=0.05] ([xshift=-5.1cm,yshift=-6.8cm]current page.north east) circle (1.4cm);
\end{tikzpicture}

\vspace*{0.9cm}
{\color{Gold}\normalsize\bfseries PROPOSITION DE SUJET DE PROJET D'EXPERTISE}\par
\vspace{1.2cm}
{\color{white}\fontsize{29}{35}\selectfont\bfseries Sécurité applicative intelligente intégrée au DevSecOps\par}
\vspace{0.6cm}
{\color{white!85}\large Agents spécialisés, tests contrôlés et intégration au cycle de développement\par}

\vspace{3.2cm}
\begin{tcolorbox}[
  enhanced,colback=white,colframe=Line,boxrule=0.6pt,arc=2mm,
  left=5mm,right=5mm,top=4mm,bottom=4mm,width=0.96\textwidth,
  borderline west={3pt}{0pt}{Teal},drop fuzzy shadow=Navy!20
]
  \begin{tabularx}{\linewidth}{@{}>{\bfseries\color{Muted}}p{3.0cm}Y@{}}
    Équipe & Badr Laklach et Taha Benahmadi \\
    Filière & GI--ILSI \\
    Destinataire & M. El Moukhtar Zemmouri, Chef de filière GI--ILSI \\
    Objet & Présentation d'une idée de projet en vue d'obtenir un accord de principe \\
  \end{tabularx}
\end{tcolorbox}

\vfill
{\color{Muted}\large Proposition académique — \today}\par
\end{titlepage}

\section{Présentation de l'idée}

\begin{summarybox}
Nous proposons d'étudier la conception d'une plateforme qui intègre des agents d'intelligence artificielle à un processus de sécurité des applications web. À partir d'une application de test autorisée, le système en cartographie les points d'entrée, propose des contrôles de sécurité adaptés, fait vérifier les résultats par des outils contrôlés et transmet un rapport à l'équipe de développement dans le cadre d'une chaîne d'intégration continue.

L'idée associe la cybersécurité applicative, l'automatisation DevOps et les pratiques DevSecOps. Les agents assistent l'analyse et la sélection des tests ; l'exécution et la confirmation des résultats restent encadrées par des mécanismes déterministes.
\end{summarybox}

\section{Contexte et objectif}

Les applications évoluent au fil de changements fréquents. Lorsqu'un contrôle de sécurité intervient tardivement, les défauts peuvent être plus difficiles à repérer et à corriger. Les outils automatisés apportent une première couverture, tandis que les agents d'IA peuvent aider à interpréter le contexte d'une application et à orienter les vérifications.

Cependant, une proposition produite par une IA ne constitue pas, à elle seule, la preuve d'une vulnérabilité. L'enjeu du projet est donc de relier l'analyse assistée par agents à des tests limités, vérifiables et intégrés au cycle de développement.

\begin{questionbox}
Comment utiliser des agents d'IA spécialisés pour orienter des contrôles de sécurité web et restituer des résultats vérifiables dans un processus DevSecOps, tout en limitant strictement les tests à un environnement autorisé ?
\end{questionbox}

L'objectif général est d'étudier et de réaliser un prototype démontrant cette logique sur une application de laboratoire isolée. Le système doit aider à repérer des problèmes de sécurité, fournir des éléments permettant de les examiner et communiquer les résultats au cours du processus de livraison logicielle.

\section{Logique générale du système}

Le fonctionnement envisagé suit un parcours simple. Une exécution est déclenchée dans un périmètre défini ; l'application de laboratoire est préparée ; le système repère ses pages et points d'entrée ; des agents spécialisés proposent les vérifications pertinentes ; un contrôle de périmètre précède toute interaction avec la cible ; enfin, les tests sont vérifiés et leurs résultats sont restitués dans le pipeline.

\begin{center}
\begin{tikzpicture}[
  node distance=5mm,
  stage/.style={
    draw=Blue!55,rounded corners=2mm,fill=PaleBlue,
    text width=3.05cm,align=center,
    inner sep=2.2mm,font=\small
  },
  guard/.style={stage,draw=Teal,line width=1pt,fill=PaleTeal},
  output/.style={stage,draw=Gold,fill=PaleGold},
  arrow/.style={-{Stealth[length=2mm]},line width=0.8pt,color=Teal}
]
  \node[stage] (a) {\textbf{1. Déclenchement}\\Un changement logiciel ou une exécution autorisée lance le contrôle.};
  \node[stage,right=of a] (b) {\textbf{2. Découverte et analyse}\\L'application isolée est cartographiée ; les agents proposent des contrôles adaptés.};
  \node[guard,right=of b] (c) {\textbf{3. Contrôle et vérification}\\Le périmètre est contrôlé avant l'exécution ; les résultats sont confrontés à des preuves.};
  \node[output,right=of c] (d) {\textbf{4. Restitution}\\Un rapport clair est transmis au pipeline et à l'équipe de développement.};
  \draw[arrow] (a.east) -- (b.west);
  \draw[arrow] (b.east) -- (c.west);
  \draw[arrow] (c.east) -- (d.west);
\end{tikzpicture}
\end{center}

Le principe important est la séparation entre \textbf{proposition} et \textbf{action} : les agents analysent et suggèrent ; un mécanisme de contrôle autorise les interactions ; des outils réalisent les tests ; les preuves, et non la seule réponse du modèle, déterminent le résultat présenté.

\section{Apport des trois domaines}

\begin{center}
\begin{tabularx}{\textwidth}{@{}>{\bfseries\color{Blue}}YYY@{}}
\hline
Cybersécurité & DevOps & DevSecOps \\
\hline
Comprendre et contrôler les vulnérabilités d'applications web, puis étayer les constats par des preuves. & Préparer un environnement reproductible et automatiser le lancement des contrôles et la collecte des résultats. & Insérer les vérifications de sécurité dans le cycle de livraison et donner aux développeurs un retour exploitable. \\
\hline
\end{tabularx}
\end{center}

Ces dimensions donnent au sujet une portée transversale : il traite à la fois de la protection des applications, de l'automatisation des environnements et de la prise en compte de la sécurité pendant le développement.

\section{Cadre envisagé}

Les essais porteraient exclusivement sur une application de laboratoire volontairement vulnérable et maîtrisée par l'équipe. Le projet ne viserait pas des applications publiques, des services de tiers ou des environnements de production.

Les familles de vulnérabilités et le niveau d'automatisation seraient définis après l'accord de principe, avec l'encadrement, en fonction de la durée du projet et des objectifs pédagogiques. Des exemples possibles sont les défauts de contrôle d'accès ou les injections dans les applications web ; ils ne constituent pas encore un périmètre définitif.

La plateforme aurait pour finalité d'assister l'évaluation et le retour aux développeurs. Toute décision importante resterait fondée sur des règles explicites et sur des résultats vérifiables.

\section{Demande d'accord de principe}

Ce sujet réunit des compétences de la filière GI--ILSI en cybersécurité, DevOps, DevSecOps, intelligence artificielle et génie logiciel. Il propose d'étudier une question actuelle tout en gardant une finalité pratique : intégrer des contrôles de sécurité web vérifiables au cycle de développement.

Nous souhaitons soumettre cette idée à l'appréciation de M. El Moukhtar Zemmouri, Chef de filière GI--ILSI, et solliciter un accord de principe pour poursuivre le cadrage du sujet. Le nom éventuel du projet, son périmètre précis et ses choix techniques resteront à définir avec l'encadrement après cette première validation.

\vfill
\begin{center}
\begin{tcolorbox}[
  colback=PaleBlue,colframe=Blue!30,boxrule=0.5pt,arc=2mm,
  width=0.94\textwidth,left=3mm,right=3mm,top=2mm,bottom=2mm
]
\centering\small\color{Muted}
Proposition de principe — GI--ILSI\\
Badr Laklach et Taha Benahmadi
\end{tcolorbox}
\end{center}

\end{document}
